
Neural networks trained with empirical risk minimization exploit spurious correlations, degrading performance on minority groups where training-time patterns fail to hold. A common explanation posits that spurious features produce stronger gradient signals, outcompeting invariant features during learning. We test this hypothesis directly by measuring gradient norms for spurious-aligned versus minority samples during training on Waterbirds. Our experiments confirm spurious feature dominance from epoch 0 (GradCAM attribution ratio 1.35) but falsify the gradient competition mechanism: minority samples produce 5.$7\times$ higher gradient norms than spurious-aligned samples, inverting the hypothesized direction. Across three random seeds, the observed gradient ratio is 0.18, not the >1.5 predicted by gradient competition. We propose that spurious dominance arises from convergence speed---simpler patterns reach low loss faster---rather than gradient magnitude. This reframing suggests that robustness interventions should target loss landscape geometry or convergence timing rather than gradient rebalancing.
